% Manual do Usuário — Planck Automation
% Compilar com: pdflatex manual_do_usuario.tex  (rodar duas vezes pelo sumário)
\documentclass[11pt,a4paper]{article}

\usepackage[utf8]{inputenc}
\usepackage[T1]{fontenc}
\usepackage[brazil]{babel}
\usepackage[margin=2.5cm]{geometry}
\usepackage{amsmath}
\usepackage{booktabs}
\usepackage{xcolor}
\usepackage{enumitem}
\usepackage{longtable}
\usepackage[colorlinks=true,linkcolor=azulesc,urlcolor=azulesc]{hyperref}

\definecolor{azulesc}{HTML}{0D47A1}
\definecolor{cinzaclaro}{HTML}{F0F0F0}
\definecolor{ambar}{HTML}{8D6E00}

\newcommand{\aba}[1]{\textsf{\textbf{#1}}}
\newcommand{\botao}[1]{\textsf{[#1]}}
\newcommand{\campo}[1]{\textsf{\textit{#1}}}

\newenvironment{aviso}
  {\par\medskip\noindent\begin{tabular}{|p{0.96\linewidth}|}\hline\rule{0pt}{3ex}}
  {\\[1ex]\hline\end{tabular}\par\medskip}

\title{\textbf{Planck Automation}\\[0.3em]
\large Determinação da constante de Planck por radiação de corpo negro\\[0.6em]
\Large Manual do Usuário}
\author{Documentação do software de bancada}
\date{\today}

\begin{document}
\maketitle
\thispagestyle{empty}

\begin{abstract}
\noindent
Este manual descreve como operar o software \emph{Planck Automation}: instalação,
ligação aos instrumentos, execução de uma simulação e de um experimento real,
leitura dos resultados e geração de relatórios. Para entender \emph{como} o
software funciona por dentro --- a arquitetura, o modelo físico e a lógica de
cálculo --- consulte o documento companheiro \textbf{Manual Técnico}.
\end{abstract}

\tableofcontents
\newpage

% =============================================================
\section{O que é este software}

O \textbf{Planck Automation} automatiza um experimento didático de física
moderna: determinar a constante de Planck $h$ medindo a radiação emitida por um
filamento de tungstênio aquecido.

O princípio, retirado do artigo de Cavalcante \& Haag (2005), é usar um
\textbf{LED como sensor espectral seletivo}. Um LED só responde a uma faixa
estreita de comprimentos de onda em torno de um $\lambda$ característico. Se o
filamento for aquecido a temperaturas crescentes e a fotocorrente do LED for
medida a cada temperatura, a relação entre as duas grandezas segue a
aproximação de Wien para a radiação de corpo negro. Linearizando essa relação,
a inclinação da reta resultante contém $h$.

O software executa esse procedimento de ponta a ponta:

\begin{enumerate}[nosep]
  \item aplica uma sequência crescente de tensões ao filamento, usando uma
        fonte programável;
  \item aguarda a estabilização térmica em cada ponto;
  \item lê a corrente do filamento (na fonte) e a fotocorrente do LED (no
        multímetro);
  \item calcula a resistência $R = V/i$ e, a partir dela, a temperatura do
        filamento;
  \item grava tudo em um arquivo CSV, ponto a ponto, enquanto coleta;
  \item ao final, faz a regressão linear e apresenta $h$, o erro relativo em
        relação ao valor de referência e o coeficiente $R^2$.
\end{enumerate}

\subsection{Instrumentos}

\begin{center}
\begin{tabular}{@{}lll@{}}
\toprule
\textbf{Instrumento} & \textbf{Modelo} & \textbf{Papel} \\
\midrule
Fonte de alimentação & Tektronix PWS4323 & Aquece o filamento (0--32~V, 0--3~A) \\
Multímetro & Tektronix DMM4050 & Lê a fotocorrente do LED \\
\bottomrule
\end{tabular}
\end{center}

A fonte conecta por \textbf{USB}; o multímetro, por \textbf{rede TCP/IP}. Os
dois falam o protocolo SCPI através da camada VISA.

% =============================================================
\section{Instalação e execução}

\subsection{Requisitos}

\begin{itemize}[nosep]
  \item Windows com Python 3.10 ou superior (o ambiente de referência usa 3.13);
  \item driver VISA instalado (NI-VISA ou TekVISA) --- necessário apenas para
        operar a bancada real;
  \item as bibliotecas listadas em \texttt{Software/requirements.txt}:
        PySide6, PyVISA, numpy, pandas, pyqtgraph e reportlab.
\end{itemize}

\subsection{Preparando o ambiente}

Na primeira vez, dentro da pasta \texttt{Software}:

\begin{quote}\ttfamily
python -m venv .venv\\
.venv$\backslash$Scripts$\backslash$activate\\
pip install -r requirements.txt
\end{quote}

\subsection{Executando}

\begin{quote}\ttfamily
cd Software\\
python main.py
\end{quote}

A janela abre com \textbf{quatro abas} no topo. O tema é escuro por padrão.

\begin{aviso}
\textbf{Importante:} execute sempre a partir da pasta \texttt{Software}. Os
arquivos de coleta são gravados em \texttt{data\_backup/} relativo ao diretório
de trabalho.
\end{aviso}

% =============================================================
\section{Visão geral da janela}

\begin{center}
\begin{tabular}{@{}p{4.2cm}p{10.5cm}@{}}
\toprule
\textbf{Aba} & \textbf{Para que serve} \\
\midrule
\aba{Hardware e Ligações} & Encontrar e validar os instrumentos; ligar ou
desligar o modo demonstração. \\[0.4em]
\aba{Simulação} & Gerar dados sintéticos e ver toda a cadeia de cálculo
funcionando, sem instrumentos. \\[0.4em]
\aba{Experimento Real} & Executar a varredura na bancada e coletar medidas
de verdade. \\[0.4em]
\aba{Referências} & Consultar o artigo de base e os manuais dos instrumentos,
com o caminho oficial de cada um. \\
\bottomrule
\end{tabular}
\end{center}

% =============================================================
\section{Aba Hardware e Ligações}

É por aqui que se começa. A aba tem duas caixas de seleção (uma para cada
instrumento), dois botões e um interruptor de modo.

\subsection{Encontrando os instrumentos}

\begin{description}[style=nextline,leftmargin=0.6cm]
  \item[\botao{Procurar Instrumentos}]
    Varre o barramento VISA à procura da fonte (por USB) e testa o endereço de
    rede do multímetro. O resultado preenche as duas caixas de seleção.
    Enquanto uma varredura está em curso, cliques adicionais no botão são
    ignorados --- não é necessário clicar repetidamente.

  \item[\botao{Verificar Ligações}]
    Pergunta a identificação (\texttt{*IDN?}) a cada instrumento selecionado e
    mostra o resultado no indicador ao lado da caixa.
\end{description}

\subsection{Os indicadores de estado}

\begin{center}
\begin{tabular}{@{}cl@{}}
\toprule
\textbf{Indicador} & \textbf{Significado} \\
\midrule
Círculo branco & Ainda não verificado \\
Círculo verde & Instrumento respondeu; ligação válida \\
Círculo amarelo & Instrumento \emph{simulado} (modo demonstração) \\
Círculo vermelho & Falha na ligação \\
\bottomrule
\end{tabular}
\end{center}

Passar o mouse sobre o indicador mostra a resposta completa do instrumento (ou
a mensagem de erro, se falhou).

\subsection{As caixas de seleção são editáveis}

Se a varredura não encontrar o multímetro --- por exemplo, porque o endereço de
rede mudou --- é possível digitar o endereço diretamente na caixa, no formato:

\begin{quote}\ttfamily
TCPIP::192.168.1.107::3490::SOCKET
\end{quote}

e clicar em \botao{Verificar Ligações}. O último par nome/endereço validado com
sucesso é lembrado para a próxima vez que o software abrir.

\subsection{Modo demonstração}
\label{sec:demo}

O interruptor \campo{Modo demonstração --- bancada simulada, sem hardware}
substitui os dois instrumentos por versões virtuais. Com ele ligado:

\begin{itemize}[nosep]
  \item nenhum comando é enviado pela rede ou pela USB;
  \item as caixas passam a mostrar \emph{PWS4323 (simulado)} e
        \emph{DMM4050 (simulado)};
  \item os indicadores ficam amarelos;
  \item a aba \aba{Experimento Real} funciona normalmente, mas os dados vêm de
        um filamento e um LED virtuais.
\end{itemize}

O filamento simulado tem valores fixos, informados na própria tela:
$R_0 = 1{,}2~\Omega$, $\alpha = 5{,}23 \times 10^{-3}$,
$\beta = 7{,}0 \times 10^{-7}$ e um LED de $\lambda = 590$~nm. \textbf{Use esses
mesmos valores nos parâmetros do experimento}, senão a temperatura calculada
não corresponderá à temperatura real do filamento virtual.

\begin{aviso}
Coletas em modo demonstração são gravadas com o prefixo
\texttt{demo\_planck\_} em vez de \texttt{exp\_planck\_}, para que dados
simulados nunca se misturem ao acervo de medidas reais.
\end{aviso}

O estado do interruptor é lembrado entre sessões. Ao desligá-lo, o software
recupera o último instrumento real validado.

% =============================================================
\section{Aba Simulação}

Gera dados sintéticos e roda a mesma cadeia de cálculo do experimento real.
Serve para conhecer o fluxo, testar parâmetros e demonstrar o método sem
ocupar a bancada.

\subsection{Parâmetros}

A aba se divide em duas sub-abas: \campo{1. Parâmetros e Varredura} e
\campo{2. Coleta e Resultados}.

\begin{center}
\begin{longtable}{@{}p{4.6cm}p{2.2cm}p{7.5cm}@{}}
\toprule
\textbf{Campo} & \textbf{Padrão} & \textbf{O que é} \\
\midrule
\endhead
Resistência a frio $R_0$ & 1,2~$\Omega$ &
Resistência do filamento usada como referência no modelo $R(T)$. \\[0.3em]
Coef. linear $\alpha$ & 5,23e-3~K$^{-1}$ &
Primeiro coeficiente de temperatura da resistência do tungstênio. \\[0.3em]
Coef. quadrático $\beta$ & 7,0e-7~K$^{-2}$ &
Segundo coeficiente, que dá a curvatura de $R(T)$. \\[0.3em]
Comprimento de onda $\lambda$ & 590~nm &
Comprimento de onda característico do LED usado como sensor. \\[0.3em]
Fator de ruído & 0,05 &
Amplitude do ruído aleatório somado à fotocorrente, como fração do sinal. \\[0.3em]
Tensão inicial & 0,5~V & Primeiro ponto da varredura. \\[0.3em]
Tensão final & 12,0~V & Último ponto da varredura. \\[0.3em]
Passo de tensão & 0,1~V & Incremento entre pontos. \\[0.3em]
Intervalo de captura & 50~ms &
Pausa entre pontos, apenas para a animação dos gráficos. \\
\bottomrule
\end{longtable}
\end{center}

\subsection{Executando}

Clique em \botao{Iniciar Coleta Simulada}. O software muda automaticamente para
a sub-aba de resultados e começa a desenhar os pontos. \botao{Parar Coleta}
interrompe a qualquer momento.

Dois gráficos são atualizados em tempo real:

\begin{itemize}[nosep]
  \item \textbf{Monitor de Fotocorrente}: a fotocorrente medida contra a tensão
        aplicada --- a curva crua, fortemente não linear;
  \item \textbf{Linearização $\ln(I)$ vs $1/T$}: os mesmos dados na forma
        linearizada. É a inclinação \emph{desta} reta que fornece $h$.
\end{itemize}

Ao terminar, a reta de ajuste aparece tracejada sobre os pontos, e os dois
painéis superiores mostram o valor de $h$ obtido, o erro relativo e o $R^2$.

% =============================================================
\section{Aba Experimento Real}

\begin{aviso}
\textbf{Segurança.} O filamento chega a mais de 2500~K. Antes de iniciar,
confirme que a montagem está correta e que o limite de corrente da fonte está
adequado ao filamento em uso. O software zera a tensão e desliga a saída da
fonte em qualquer caminho de erro ou parada, inclusive se a janela for fechada
durante a coleta --- mas isso não substitui a conferência da montagem.
\end{aviso}

\subsection{Antes de começar}

\begin{enumerate}[nosep]
  \item Ligue os dois instrumentos e confirme na aba \aba{Hardware e Ligações}
        que ambos aparecem em verde.
  \item Meça a resistência a frio do filamento e anote o valor.
  \item Confirme o comprimento de onda do LED que será usado como sensor.
\end{enumerate}

\subsection{Parâmetros}

A sub-aba \campo{1. Parâmetros Reais} tem dois grupos.

\paragraph{Calibração Física do Filamento.}
Os mesmos quatro campos da simulação: $R_0$, $\alpha$, $\beta$ e $\lambda$. São
eles que convertem a resistência medida em temperatura e a inclinação da reta
em $h$.

\paragraph{Varredura SCPI.}

\begin{center}
\begin{tabular}{@{}p{4.6cm}p{2.2cm}p{7.5cm}@{}}
\toprule
\textbf{Campo} & \textbf{Padrão} & \textbf{O que é} \\
\midrule
Tensão inicial & 1,0~V & Primeira tensão aplicada ao filamento. \\[0.3em]
Tensão final & 10,0~V & Última tensão da varredura. \\[0.3em]
Passo de tensão & 0,5~V & Incremento entre pontos. \\[0.3em]
Estabilização térmica & 3000~ms &
\textbf{Tempo de espera em cada ponto antes de medir.} O filamento leva
segundos para atingir o equilíbrio térmico; medir antes disso produz
temperaturas erradas. \\
\bottomrule
\end{tabular}
\end{center}

O número de pontos é determinado pelos três primeiros campos. Uma varredura de
1,0~V a 10,0~V com passo 0,5~V produz 19 pontos; com 3~s de espera cada, a
coleta leva cerca de um minuto.

\subsection{Executando}

\botao{Iniciar Experimento Físico} liga a saída da fonte e começa a varredura.
A tela muda para a sub-aba \campo{2. Coleta na Bancada}, que mostra:

\begin{itemize}[nosep]
  \item uma \textbf{barra de progresso} com o ponto atual sobre o total;
  \item três \textbf{gráficos ao vivo}: temperatura contra tensão, fotocorrente
        contra tensão, e a linearização $\ln(I)$ vs $1/T$;
  \item um \textbf{registro em tempo real}, à direita, com uma linha por ponto
        medido, no formato \texttt{tensão | corrente | temperatura |
        fotocorrente};
  \item o nome do arquivo em que os dados estão sendo gravados.
\end{itemize}

\botao{Parar e Processar Experimento} interrompe a varredura, desliga a fonte
com segurança e \textbf{ainda assim calcula} o resultado com os pontos já
coletados --- desde que haja pelo menos três.

\subsection{Onde ficam os dados}

Cada execução cria um arquivo em \texttt{Software/data\_backup/}:

\begin{quote}\ttfamily
data\_backup/exp\_planck\_AAAAMMDD\_HHMMSS.csv
\end{quote}

O arquivo é escrito \textbf{ponto a ponto, durante a coleta}, não ao final.
Se houver queda de energia ou falha de comunicação no meio do experimento, tudo
o que já foi medido está salvo. As colunas são:

\begin{center}
\begin{tabular}{@{}ll@{}}
\toprule
\textbf{Coluna} & \textbf{Conteúdo} \\
\midrule
\texttt{Tensao\_Fonte\_V} & Tensão programada na fonte, em volts \\
\texttt{Corrente\_Filamento\_A} & Corrente lida na fonte, em ampères \\
\texttt{Resistencia\_Ohms} & $R = V/i$, em ohms \\
\texttt{Temperatura\_K} & Temperatura calculada do filamento, em kelvin \\
\texttt{Fotocorrente\_A} & Corrente do LED lida no multímetro, em ampères \\
\bottomrule
\end{tabular}
\end{center}

% =============================================================
\section{Lendo os resultados}
\label{sec:resultados}

Ao final, o software mostra três números:

\begin{description}[style=nextline,leftmargin=0.6cm]
  \item[$h$ experimental]
    O valor obtido, em J$\cdot$s. O valor de referência (CODATA) é
    $6{,}62607015 \times 10^{-34}$~J$\cdot$s.
  \item[Erro relativo]
    A diferença percentual entre o valor obtido e o de referência.
  \item[$R^2$]
    Quão bem os pontos se ajustam a uma reta na forma linearizada. Um $R^2$
    próximo de 1 indica que a relação linear se sustentou; um $R^2$ baixo
    indica que a regressão pegou pontos que não seguem o modelo.
\end{description}

\subsection{Por que a faixa de varredura importa tanto}
\label{sec:faixa}

Esta é a característica mais importante de operação do experimento, e vale
entendê-la antes de interpretar qualquer resultado.

A aproximação de Wien só descreve bem a fotocorrente quando o filamento está
\textbf{quente}. Em temperaturas baixas, a radiação emitida no comprimento de
onda do LED é ordens de grandeza menor que o \textbf{ruído de fundo do
multímetro}. O que se lê nesses pontos não é sinal: é o piso do instrumento.

Na prática, com o multímetro na faixa de 100~$\mu$A, o fundo é da ordem de
alguns nanoampères. Uma varredura completa a partir de 0~V produz, portanto,
dois regimes bem distintos:

\begin{center}
\begin{tabular}{@{}lll@{}}
\toprule
\textbf{Região} & \textbf{Fotocorrente típica} & \textbf{Vale como medida?} \\
\midrule
Abaixo de $\approx 1800$~K & unidades de nA & Não --- é o piso do multímetro \\
Acima de $\approx 1800$~K & dezenas a centenas de nA & Sim \\
\bottomrule
\end{tabular}
\end{center}

Incluir os pontos frios na regressão não apenas não ajuda: \textbf{estraga o
resultado}, porque eles entram como se fossem medidas legítimas e puxam a
reta. O sintoma característico é um $R^2$ baixo acompanhado de um $h$ muito
distante do esperado.

\begin{aviso}
\textbf{Recomendação prática.} Você pode varrer a faixa inteira que quiser ---
inclusive de 0 a 12~V --- para registrar o comportamento completo do filamento.
Mas, ao avaliar o valor de $h$, confira o $R^2$: se ele estiver baixo, refaça a
leitura considerando apenas a região quente. O CSV guarda todos os pontos, então
nada se perde.
\end{aviso}

% =============================================================
\section{Exportando o relatório PDF}

Depois que uma coleta (simulada ou real) termina, o botão
\botao{Exportar Relatório PDF} é habilitado. Ele abre primeiro uma janela de
metadados:

\begin{itemize}[nosep]
  \item \campo{Autor / Pesquisador};
  \item \campo{Local da Coleta};
  \item \campo{Notas Adicionais} --- espaço livre para registrar anomalias no
        equipamento, temperatura ambiente, condições da sala.
\end{itemize}

Em seguida, um seletor de arquivos pergunta onde salvar. O PDF gerado contém:

\begin{enumerate}[nosep]
  \item cabeçalho com pesquisador, local e data/hora do registro;
  \item tabela com todos os parâmetros usados na coleta;
  \item tabela de resultados ($h$ de referência, $h$ experimental, erro
        relativo e $R^2$);
  \item uma imagem dos gráficos, capturada da tela no momento da exportação;
  \item as observações escritas na janela de metadados.
\end{enumerate}

% =============================================================
\section{Aba Referências}

Lista os quatro documentos em que o software se baseia: o artigo de Cavalcante
\& Haag, os datasheets do multímetro e da fonte, e o manual do TekVISA.

Para cada documento, a ficha traz a autoria, onde foi publicado, o
identificador (DOI ou número de peça), a condição de acesso, e uma lista do que
especificamente o software usa daquele documento --- quais equações, quais
tabelas de exatidão.

Três botões acompanham cada ficha:

\begin{description}[style=nextline,leftmargin=0.6cm,nosep]
  \item[\botao{Abrir cópia local}] abre o PDF guardado na pasta \texttt{Docs/}.
  \item[\botao{Abrir documento oficial}] abre o link direto do documento no
        site do editor ou do fabricante.
  \item[\botao{Página do editor/fabricante}] abre a página que hospeda o
        documento, útil para verificar se houve revisão.
\end{description}

% =============================================================
\section{Solução de problemas}

\begin{description}[style=nextline,leftmargin=0.6cm]

  \item[A varredura não encontra a fonte]
    Confirme que o cabo USB está conectado e que o driver VISA está instalado.
    A fonte é identificada pelo código de fabricante Tektronix no endereço USB.

  \item[A varredura não encontra o multímetro]
    O endereço de rede do multímetro é procurado num endereço padrão. Se o
    aparelho estiver em outro IP, digite o endereço completo diretamente na
    caixa de seleção e clique em \botao{Verificar Ligações}.

  \item[Aparece ``valide a ligação aos instrumentos primeiro'']
    O experimento real exige que os dois instrumentos tenham sido validados com
    sucesso ao menos uma vez. Vá à aba \aba{Hardware e Ligações} e use
    \botao{Verificar Ligações}. Alternativamente, ligue o modo demonstração
    (seção~\ref{sec:demo}) para rodar sem hardware.

  \item[O multímetro fica travado em modo remoto]
    Não deve mais acontecer: ao encerrar a comunicação, o software devolve o
    controle ao painel frontal. Se ocorrer, o botão \textsc{local} do painel do
    aparelho resolve.

  \item[O resultado de $h$ saiu muito distante do esperado]
    Verifique o $R^2$. Se estiver baixo, provavelmente a varredura incluiu
    pontos frios demais --- veja a seção~\ref{sec:faixa}. Se o $R^2$ estiver
    alto mas $h$ ainda estiver deslocado, o problema tende a ser de calibração:
    o valor de $R_0$ e o comprimento de onda do LED são os parâmetros que mais
    influenciam o resultado.

  \item[Temperaturas absurdas em tensão baixa]
    Em tensões próximas de zero, a resistência medida pode ficar abaixo de
    $R_0$, e o modelo devolve temperaturas sem sentido físico (por exemplo,
    77~K). Esses pontos devem ser desconsiderados; não indicam falha do
    equipamento.

  \item[A janela fechou durante a coleta]
    O software interrompe a varredura e desliga a fonte com segurança antes de
    encerrar. Os dados já coletados estão no CSV.

\end{description}

\vfill
\begin{center}
\rule{0.5\linewidth}{0.4pt}\\[0.5em]
\small Documento companheiro: \textbf{Manual Técnico}, que descreve a
arquitetura interna e o modelo físico.
\end{center}

\end{document}
